\section{Adaptive Thresholding Approaches} \label{sec:ATA}

This section describes the denoising and binarization algorithms investigated in this work.
All implemented approaches follow their original mathematical formulations; however, the underlying implementations are adapted to improve efficiency and to accommodate the processing of color image data.

% First pre-process (colored image (r,g,b) to gray scale value)
As an initial preprocessing step, all input color images are converted to grayscale. 
This allows subsequent denoising and binarization algorithms to operate on a single intensity channel, simplifying the computation of local statistics while preserving the relevant structural information of the image.

The grayscale intensity is computed from the RGB color channels using a linear luminance approximation. 
This approach combines the red, green, and blue components using fixed weighting factors, reflecting their relative contribution to perceived brightness.
Following conversion, the resulting grayscale image is normalized to an 8-bit intensity range to ensure consistent scaling across different inputs and formats.

% plug in respective formulas
This preprocessing step is applied uniformly to all evaluated approaches and image formats and does not alter the underlying denoising or thresholding formulations.
By standardizing the input representation, it enables a fair comparison of the considered algorithms and isolates the effects of the thresholding techniques and implementation strategies discussed in the subsequent sections.

% Next sections: Explain shortly which approaches we looked at:
% - Baseline:
%     - Niblack
%     - Sauvola
%     - Nick
% - SOTA:
%     - binarization
%     - integral